\subsection{A multisignature Force and the retained anchor}
\label{sec:force-example}
Let the region order be A, B, C. A's session $(7,20)$ produces
DATA 101--110. Its terminal receipt lets B start $(7,21)$ before
all A DATA is globally accepted. Suppose B has produced through 115,
while $G$ ends at 108. The public export session is still A's:
the next required DATA is 109.

The committee submits independent OFF(A) and MSIG targeting version 7,
both for activation threshold 120. Both satisfy advance announcement.
Before the first finalized Espresso block reaching 120, valid old
output may still extend $G$.

If $G$ still ends at 108 at that boundary, replay applies OFF(A),
checks the target version and closes version 7 before payload.
It retains 108. Old A:109--110 and B:111--115 can no longer append.
Their receipts remain historical evidence of regional commitment.

If B is eligible, the forced target is $(8,21)$ because A already
contributed positive accepted DATA. This new engine enters
$(8,21,0,0)$ from 108 and may propose DATA 109 after Ready and the
original ANCHOR bootstrap. It does not wait for Espresso results to
be registered on L1. New-version DATA is accepted only in a later
Espresso block.

If another valid MSIG already changed version 7 to 8 before 120,
the later instruction targeting 7 is stale. Independent OFF(A) still
takes effect. If A completed its export and B became responsible through
regular rotation, the version remains 7: the instruction still executes,
retains the newer $G$ and selects after the actual export region B.

With quota ten, old B $(7,21)$ starts after 110 and ends at 120.
New B $(8,21)$ starts after 108 and ends at 118. They use the same
immutable authority index but different session engines and anchors.

\begin{figure}[p]\centering
% Explicit MSIG control and the existing regional bootstrap.
% beta is the numeric branch version, and (beta,r) identifies a session.
% Vertical spacing illustrates event order, not elapsed time.
\begin{tikzpicture}[
  x=1cm,y=-1cm,font=\fontsize{8.3}{9.8}\selectfont,
  party/.style={draw,rounded corners=2pt,fill=white,align=center,
    inner sep=3pt,text width=1.22cm,minimum height=.66cm},
  lane/.style={draw=black!28,densely dashed},
  msg/.style={-{Stealth[length=1.6mm]},semithick,draw=rrBlue},
  read/.style={msg,densely dashed},
  force/.style={msg,draw=rrRed},
  tag/.style={fill=white,inner sep=1.5pt,font=\fontsize{7.7}{9}\selectfont},
  note/.style={anchor=north west,align=flush left,text width=5.55cm,
    inner sep=0pt,font=\fontsize{8.3}{9.8}\selectfont},
  phase/.style={anchor=west,fill=white,inner sep=2pt,
    font=\fontsize{9}{10.5}\selectfont\bfseries},
  dot/.style={circle,fill=black,inner sep=1.1pt},
  wait/.style={line width=2.5pt,draw=rrTeal!75,
    dash pattern=on 2.5pt off 1.5pt},
  rulebox/.style={draw=rrRed!55,rounded corners=2pt,fill=rrRed!3,
    align=center,text width=7.9cm,inner sep=5pt}
]
  \def\xC{.75}\def\xI{2.3}\def\xE{3.85}
  \def\xN{5.4}\def\xL{6.95}\def\xB{8.5}\def\xn{9.45}
  \path[use as bounding box] (-.05,-.55) rectangle (15.05,16.05);
  \node[party] at (\xC,0) {Multisig};
  \node[party] at (\xI,0) {L1};
  \node[party] at (\xE,0) {Espresso};
  \node[party] at (\xN,0) {Node /\\certifier};
  \node[party] at (\xL,0) {B: initial\\leader};
  \node[party] at (\xB,0) {B: other\\replicas};
  \foreach \x in {\xC,\xI,\xE,\xN,\xL,\xB}
    \draw[lane] (\x,.52)--(\x,15.65);

  \node[phase] at (0,.88) {1.\ Announce MSIG Force on L1};
  \draw[force] (\xC,1.55)--(\xI,1.85);
  \node[tag,anchor=south] at (1.52,1.48) {submit $F$ at $s$};
  \node[note] at (\xn,1.30) {$F$ names target version $\beta$ and future
    time $\theta_a$. It applies only if $\beta$ is still current.
    Independent STATUS keeps its own effect if the Force is stale.};
  \node[dot] at (\xI,2.65) {};
  \draw[read,draw=rrRed] (\xI,2.65)--node[tag,above,sloped]
    {finalized control + complete history}(\xN,3.30);
  \node[dot] at (\xN,3.30) {};
  \node[note] at (\xn,2.85) {By $s+T_{\rm L1}$, all correct participants
    can verify the finalized control. $T_{\rm L1}$ includes dissemination
    and verification. This precedes the Espresso crossing.};
  \node[tag,align=center] at (5.35,4.15)
    {While version $\beta$ remains current: output stays eligible;\\
     normal regional handoffs may continue.};
  \node[phase] at (0,5.25) {2.\ Finalize the Espresso boundary; derive one Force};
  \draw[force] (\xE,5.95)--(\xN,6.25);
  \node[tag,anchor=south] at (4.62,5.90) {$J_{q^*}$ finalized};
  \node[note] at (\xn,5.75) {For MSIG Force,
    $q^*=\min\{q:\theta_q\ge\theta_a\}$.
    Interpret the complete history through $J_{q^*-1}$ first.};
  \node[rulebox] at (4.65,7.30) {At the start of $J_{q^*}$:\\[2pt]
    due independent STATUS $\;\to\;$ first matching MSIG;\\
    no matching instruction means no cross-region Force;\\
    only then process EXPORT records.};
  \node[note] at (\xn,7.05) {Health detection belongs to multisig policy. A report or missing
    DATA cannot execute Force. No Ready, output or minimum-gap check
    limits a valid instruction.};
  \node[dot,fill=rrRed] at (\xN,8.55) {};
  \node[tag,anchor=east,text=rrRed] at (\xN-.12,8.55)
    {anchor $:=G$ before payload};
  \node[note] at (\xn,8.70) {Close version $\beta$ for appends; keep $G$.
    Advance to $\beta+1$. Use NextLive to select eligible B, or wait.
    Old engines close separately.};
  \draw[read] (\xN,9.60)--node[tag,pos=.55,above,sloped]
    {complete boundary}(\xL,10.00);
  \draw[read] (\xN,9.60)--(\xB,10.15);
  \node[note] at (\xn,10.15) {Duty $(\beta_b,r_b)$ has accepted count $a_b$:
    $r'=r_b$ if $a_b=0$, else $r'=r_b+1$.
    Session $(\beta+1,r')$ can append only in a later Espresso block.};

  \node[phase] at (0,11.10) {3.\ B uses the existing ANCHOR bootstrap};
  \node[dot] at (\xL,11.65) {};
  \node[dot] at (\xB,11.85) {};
  \draw[msg] (\xB,12.20)--node[tag,above,sloped] {STATUS}(\xL,12.45);
  \draw[wait] (\xL+.18,11.65)--(\xL+.18,13.15);
  \node[tag,anchor=east] at (\xL-.12,12.85) {wait $2\Delta_{\rm B}$};
  \node[note] at (\xn,11.75) {B verifies the boundary, anchor and state.
    At Ready, enter $(\beta+1,r',0,0)$. The leader waits
    $2\Delta_{\rm B}$, then sends NEWVIEW.};
  \draw[msg] (\xL,13.20)--node[tag,above,sloped] {NEWVIEW}(\xB,13.45);
  \draw[msg] (\xB,14.00)--node[tag,above,sloped] {ANCHOR votes}(\xL,14.25);
  \node[dot] at (\xL,14.25) {};
  \node[note] at (\xn,13.65) {An ANCHOR QC with weight $>F_{\rm B}$
    permits DATA after $G$. ANCHOR votes start no commit timer.
    DATA and DRAIN retain their rules.};
  \draw[msg] (\xL,15.00)--node[tag,above,sloped]
    {DATA at $h(G)+1$}(\xB,15.30);
  \node[note] at (\xn,15.20) {The forced handoff preserves $G$;
    it does not wait for $G$ to be registered on L1.};
\end{tikzpicture}

\caption{Explicit Force: L1 supplies the instruction, Espresso time
fixes its cutoff, and the next region inherits exactly pre-payload $G$.
There is no automatic trigger and no second Espresso-to-L1 wait.}
\label{fig:forced-cross-region}
\end{figure}
\FloatBarrier

\subsection{Concrete view updates}
\label{sec:force-view-examples}
The following examples assume complete authority and an eligible target.
Every Force has an explicit valid MSIG cause.

\paragraph{A has positive accepted output.}
A is in $(7,20,2,3)$, producing 101--110, while $G=108$.
Its export session is 20 with positive accepted count. Selecting B gives
\[
 (7,20,2,3)\xrightarrow{\mathrm{Force}}(8,21,0,0).
\]
B continues after 108. Only the new engine's $g,t$ start at zero;
the retained DATA keeps its original provenance.

\paragraph{Regular rotation has already reached B.}
A first hands off
\[
 (7,20,0,2)\xrightarrow{\mathrm{regular}}(7,21,0,0).
\]
B reaches $(7,21,0,1)$ and contributes DATA through 113 to $G$.
A Force selecting C gives
\[
 (7,21,0,1)\xrightarrow{\mathrm{Force}}(8,22,0,0).
\]
C continues with DATA 114. A scheduled MSIG targeting version 7
still matches because regular rotation did not change $\beta$.

\paragraph{Speculative production is ahead of the export session.}
A $(7,20)$ has regional DATA 101--110; B $(7,21)$ has 111--120;
C runs $(7,22,1,2)$. If $G=108$, the next required DATA is still
A's 109. A Force selecting B opens $(8,21,0,0)$, not $(8,23,0,0)$.
It closes the entire old version, including C's speculative program.
Using a lower turn index than the latest speculative one removes no
globally accepted DATA.

\paragraph{A stale instruction changes no view field.}
One MSIG targets version 7 at threshold 120. A different valid MSIG
has already changed the version to 8 at an earlier position.
Suppose the new engine is now $(8,21,1,2)$. Processing the stale
instruction leaves
\[
 (8,21,1,2)\xrightarrow{\mathrm{stale\ MSIG}}(8,21,1,2).
\]
Any independent STATUS is still evaluated on its own terms.

\subsection{Zero output, an offline region and explicit exclusion}
\label{sec:retry-example}
A's $(7,20)$ opens from DATA 100 but produces no accepted DATA.
The committee observes the failure and submits OFF(A) and MSIG.
At their valid activation, if $G$ remains 100, the target slot stays
20. NextLive can choose B, which enters $(8,20,0,0)$ with a fresh
engine. Reusing the number 20 does not mean selecting A again.

If B also adds no DATA, another explicit MSIG can choose C in
$(9,20,0,0)$. Neither B's Ready nor a first DATA is a prerequisite
to that instruction. If A's DATA 101 had instead been accepted before
the first Force, the new slot would be 21 with anchor 101.

The monitoring report itself has no effect. If no MSIG is submitted,
A's missing output does not change $\beta$. If only MSIG is submitted
and A remains the only eligible region, NextLive can select A again.
Excluding it requires explicit OFF.

\subsection{Repeated Force and an unassigned boundary}
\label{sec:repeated-force}
A valid MSIG can close a newly assigned version while its target is
still Waiting, or after Ready but before any DATA. There is no minimum
Force interval. Each new instruction still needs a matching target
and valid advance announcement.

If closing positive-output session $(7,20)$ finds no eligible region,
replay stores target slot 21 and anchor $G$ in version 8. It opens no
regional engine. A later valid MSIG can change the version to 9 while
preserving target slot 21 and the same $G$. When ON(B) eventually
enables selection, B enters $(9,21,0,0)$, not $(9,22,0,0)$.
If no ON is issued, indefinite waiting is allowed.

\subsection{Strict cutoff, equal timestamps and concurrent preparation}
\label{sec:cutoff-example}
An MSIG targets version 7 at threshold 300. B may start regular
session $(7,21)$ from A's terminal receipt before that threshold.
It may also prepare a possible Force successor, updating the candidate
state as $G$ advances. No candidate grants early entry.

At the first finalized block with timestamp at least 300, the prelude
checks the still-current version. A match closes every session of
version 7 and retains the exact preceding $G$. Old DATA in that block's
payload cannot revoke the control. Candidate state for an earlier
anchor must be updated to the actual fixed anchor before Ready.

If Espresso timestamps jump across the threshold, the first crossing
block is still the activation block. A later block with the same
timestamp does not repeat an already processed instruction. A distinct
valid instruction can execute at a later position if its own target
matches. A single prelude never chains multiple Forces.
